\documentclass[10pt,a4paper]{amsart}
\usepackage[margin=19mm]{geometry}
\usepackage{amsmath,amssymb,mathtools}
\usepackage[scr=rsfs,cal=euler]{mathalfa}
\usepackage{xfrac}
\usepackage[unicode,hidelinks,bookmarks=false]{hyperref}
\newcommand{\ie}{i.e.\ }
\newcommand{\cf}{cf.\ }
\newcommand{\resp}{respectively\ }
\newcommand{\Subsection}{\subsection}
\providecommand{\nobreakdash}{\nobreak\hbox{-}}
\setlength{\emergencystretch}{2em}
\multlinegap0pt
\usepackage{bm}
\usepackage{bbm}
\usepackage{dsfont}
\usepackage{xspace}
\usepackage{cancel}
\usepackage{mathtools}
\usepackage{amsthm}
\makeatletter
\@ifundefined{theorem}{\newtheorem{theorem}{Theorem}}{}
\@ifundefined{lemma}{\newtheorem{lemma}[theorem]{Lemma}}{}
\makeatother
\newcommand{\C}{\mathbb{C}}
\newcommand{\HS}{\operatorname{HS}}
\newcommand{\N}{\mathbb{N}}
\newcommand{\R}{\mathbb{R}}
\title[Global hypoellipticity of systems of Fourier multipliers on ]{Global hypoellipticity of systems of Fourier multipliers on compact Lie groups}
\author{Andr\'e Pedroso Kowacs}
\date{Source: arXiv:2505.16958v3 (CC BY 4.0)}
\begin{document}
\maketitle

\noindent\emph{Source and licence.} Global hypoellipticity of systems of Fourier multipliers on compact Lie groups. Version arXiv:2505.16958v3 (CC BY 4.0), distributed under the Creative Commons Attribution 4.0 International licence, as recorded in the arXiv licence metadata for this exact version and shown on the abstract page https://arxiv.org/abs/2505.16958v3. Full rights notice: licenses/arxiv-2505.16958-CC-BY-4.0.txt.\par\smallskip

\noindent\textit{Editorial scope.} Counted unit: the Theorem whose source label is \texttt{theo-vector} in the article (source0.tex lines 456--462, source label \texttt{theo-vector}), delivered together with the article's own complete original proof of it (source0.tex lines 464--524, one \texttt{proof} environment of 61 lines). No mathematics is written, repaired or re-derived: statement and proof are the article's own text line for line. Required context retained uncounted: sobolev-cup-cap (lines 386--388); lemmasingvalue (lines 525--530). Every definition, display and label of those lines is retained unchanged. Omitted from this package and not redistributed: 1--385, 389--455, 463--463, 531--1037. Nothing inside a retained excerpt is omitted or altered: every retained body is delivered exactly as the article's lines are written, so no omission receipt of any kind accompanies this package. Citations: the retained excerpts carry no citation command at all, so no bibliography entry of the article is transcribed and no reference is redistributed. Reference adaptation: none was needed and none was made; every reference inside the delivered unit resolves inside it, so no unresolved reference or citation is produced. Printed slips kept literal: no printed word, symbol, hypothesis or number of the counted statement or of its proof was altered. Layout and class: the article's own class is \texttt{amsart} with packages including inputenc, amsmath, amssymb, wrapfig, url, mathtools, graphicx, stmaryrd, amsthm, xcolor, hyperref, enumitem, comment, relsize; the wrapper is the lane's pinned \texttt{amsart} editorial wrapper at 10pt on A4, which restates the theorem-like environments the retained text needs and adds the article's own macro definitions that the retained text needs (\texttt{\textbackslash C}, \texttt{\textbackslash HS}, \texttt{\textbackslash N}, \texttt{\textbackslash R}), with extra wrapper packages (bm, bbm, dsfont, xspace, cancel, mathtools). The delivered numbering is the unit's own. Assets: no retained span contains a figure environment, a graphics-inclusion command, a PDF-page inclusion, a table or a TikZ picture; no image file is copied into this package, the package declares no asset dependency and no image file is redistributed with it. Licence: version arXiv:2505.16958v3 of this article is distributed under the Creative Commons Attribution 4.0 International licence, as recorded in the arXiv licence metadata for this exact version and shown on the abstract page https://arxiv.org/abs/2505.16958v3. A full rights notice is delivered at licenses/arxiv-2505.16958-CC-BY-4.0.txt. Characters: every retained excerpt is pure ASCII, so the delivered engine, XeLaTeX with its default Latin Modern text font, reports no missing character.
\begin{equation}\label{sobolev-cup-cap}
  \bigcap_{s\in\R} H^s(G,\C^n)=C^\infty(G,\C^n) \text{ and } \bigcup_{s\in\R} H^s(G,\C^n)=\mathcal{D}'(G,\C^n).  
\end{equation}

\begin{lemma}\label{lemmasingvalue}
Let $A\in\mathbb{C}^{r\times s}$ and $B\in\mathbb{C}^{s\times p}$, where $r,s,p\in\N$, be two complex matrices. Then 
\begin{equation*}\label{ineqsingularvalue}
        \|AB\|_{HS}\geq \lambda_{\min}[A]\|B\|_{HS}.
    \end{equation*}
\end{lemma}

\begin{theorem}\label{theo-vector}
    Let $n,m\in\N$, $G$ a compact Lie group and $D:C^\infty(G,\mathbb{C}^n)\to C^\infty(G,\mathbb{C}^m)$ a left-invariant continuous linear operator. Then $D$ is globally hypoelliptic if and only if there exist $k\in\R$ and $C>0$ such that
    \begin{equation}\label{ineq_theo}
        \lambda_{\min}[\sigma_D(\xi)]\geq C\langle\xi\rangle^k,
    \end{equation}
    for all but finitely many $[\xi]\in\widehat{G}$, where $\lambda_{\min}[\sigma_D(\xi)]$ denotes the smallest singular value of the block matrix $\sigma_D(\xi)$, for $[\xi]\in\widehat{G}$.
\end{theorem}

\begin{proof}
    First assume that inequality \eqref{ineq_theo} holds for all $[\xi]\in\widehat{G}\backslash V$, where $V\subset\widehat{G}$ is finite. Then for $u\in \mathcal{D}'(G,\C^n)$ such that $Du\in C^\infty(G,\C^m)$, and for any $s\in\mathbb{R}$, we have that
    \begin{align}\label{ineqprooftheovector1}
        +\infty>\|Du\|^2_{H^s(G,\C^m)} &= \sum_{[\xi]\in\widehat{G}}d_\xi\langle\xi\rangle^{2s} \sum_{j=1}^{m}\|\widehat{(Du)_j}(\xi)\|_{\HS}^2\notag\\
        &=\sum_{[\xi]\in\widehat{G}}d_\xi\langle\xi\rangle^{2s}\sum_{j=1}^{m}\left\|\sum_{i=1}^{n }\sigma_D(i,j,\xi)\widehat{u_i}(\xi)\right\|_{\HS}^2\notag \\
        &=\sum_{[\xi]\in\widehat{G}}d_\xi\langle\xi\rangle^{2s}\left\|\sigma_D(\xi)\widehat{u}(\xi)\right\|_{\HS}^2.
    \end{align}
        But notice that, 
        \begin{align*}
          \left\|\sigma_D(\xi)\widehat{u}(\xi)\right\|_{\HS}
           &\geq 
           \lambda_{\min}[\sigma_D(\xi)]\|\widehat{u}(\xi)\|_{\HS},
        \end{align*}
        for every $[\xi]\in\widehat{G}$, by Lemma \ref{lemmasingvalue} (proven below).
%         Indeed, if $\widehat{(\chi_\tau u)_i}(\xi)=0$, for $1\leq i\leq d_\tau$ this inequality is trivial, otherwise the inequality follows from the fact that
%         \begin{align*}
%             \left\|\sum_{i=1}^{d_{\tau}}\sigma_{D}(i,j,\xi)\widehat{(\chi_\tau u)_i}(\xi)\right\|_{\HS}^2&=\left\|\sum_{i=1}^{d_{\tau}}\sigma_{D}(i,j,\xi)\frac{\widehat{(\chi_\tau u)_i}(\xi)}{(\sum_{\ell=1}^{d_\tau}\|\widehat{(\chi_\tau u)_\ell}(\xi)\|_{\HS}^2)^{\frac{1}{2}}}\right\|_{\HS}^2{\sum_{\ell=1}^{d_\tau}\|\widehat{(\chi_\tau u)_\ell}(\xi)\|_{\HS}^2},
%         \end{align*}
%         for $1\leq j\leq d_\omega$, and
%         \begin{equation*}
%             \sum_{j=1}^{d_\omega} \left\|\sum_{i=1}^{d_{\tau}}\sigma_{D}(i,j,\xi)\frac{\widehat{(\chi_\tau u)_i}(\xi)}{(\sum_{\ell=1}^{d_\tau}\|\widehat{(\chi_\tau u)_\ell}(\xi)\|_{\HS}^2)^{\frac{1}{2}}}\right\|_{\HS}^2\geq m_{\xi}(D)^2
%         \end{equation*}
%         since 
%         \begin{equation*}
% \sum_{i=1}^{d_\tau}\left\|\frac{\widehat{(\chi_\tau u)_i}(\xi)}{(\sum_{\ell=1}^{d_\tau}\|\widehat{(\chi_\tau u)_\ell}(\xi)\|_{\HS}^2)^{\frac{1}{2}}}\right\|_{\HS}^2=\frac{1}{\sum_{\ell=1}^{d_\tau}\|\widehat{(\chi_\tau u)_\ell}(\xi)\|_{\HS}^2}\sum_{i=1}^{d_\tau}\left\|{\widehat{(\chi_\tau u)_i}(\xi)}\right\|_{\HS}^2=1.
%         \end{equation*}
        Hence \eqref{ineq_theo} and \eqref{ineqprooftheovector1} imply that
    \begin{align*}    
         +\infty>\|Du\|^2_{H^s(G,\C^m)}&\geq\sum_{[\xi]\in \widehat{G}\backslash V} d_\xi\langle\xi\rangle^{2s} \lambda_{\min}[\sigma_D(\xi)]^2\|\widehat{u}(\xi)\|_{\HS}^2\\
        & \geq C^2\sum_{[\xi]\in \widehat{G}\backslash V}d_\xi\langle\xi\rangle^{2(s+k)}\|\widehat{u}(\xi)\|_{\HS}^2.
    \end{align*}
    Since $V$ is finite, this implies that $\|u\|_{H^{s+k}(G,\C^n)}<+\infty$. Because this holds for any $s\in\mathbb{R}$, we conclude by \eqref{sobolev-cup-cap} that $u\in C^\infty(G,\C^n)$ and thus $D$ is globally hypoelliptic. 
     Suppose now that inequality \eqref{ineq_theo} does not hold. Note that by the characterization of the smallest singular value, we can write 
     \begin{align*}
         \lambda_{\min}[\sigma_D(\xi)]&=\min\left\{ \|\sigma_D(\xi)v(\xi)\|_2:v(\xi)\in (\C^{d_\xi\times 1})^{n\times 1},\,\|v(\xi)\|_2=1 \right\}\\
         &=\min\left\{\left(\sum_{j=1}^m\left\|\sum_{i=1}^n\sigma_D(i,j,\xi)v(i,\xi)\right\|_2^2\right)^{\frac{1}{2}}:v(i,\xi)\in \C^{d_\xi\times 1},\,\sum_{i=1}^n\|v(i,\xi)\|_2^2=1 \right\}.
     \end{align*}
     So, for each $\ell\in\mathbb{N}$, there exist distinct $[\xi_\ell]\in\widehat{G}$ and $v(i,\xi_\ell) \in \C^{d_\xi\times 1},\,1\leq i\leq n$ satisfying $\sum_{i=1}^n\|v(i,\xi)\|_2^2=1$  such that
    \begin{equation*}
\sum_{j=1}^m\left\|\sum_{i=1}^n\sigma_D(i,j,\xi)v(i,\xi)\right\|_2^2< \langle \xi_\ell\rangle^{-\ell},
    \end{equation*}
    for every $\ell\in \N$. Let $u\in \mathcal{D}'(G,\C^n)$ be defined by the Fourier coefficients
    \begin{align*}
        \widehat{u_i}(\xi) =\begin{cases}
            \begin{pmatrix}
                \lvert&0&\cdots\\
                v(i,\xi_\ell)&0&\cdots\\
                \lvert&0&\cdots
            \end{pmatrix}_{d_\xi\times d_\xi}&\text{if } \xi=\xi_\ell,\ \ell\in\N, \text{ and }\,1\leq i\leq n,\\
            0_{d_\xi\times d_\xi}&\text{otherwise.}
        \end{cases} 
    \end{align*}
     Then $u\in\mathcal{D}'(G,\C^n)\backslash C^\infty(G,\C^n)$, as $\|\widehat{u}(\xi_\ell)\|_{\HS}^2=\sum_{i=1}^n\|v(i,\xi)\|_2^2=1$, for all $\ell\in\mathbb{N}$, and $\widehat{u}(\xi)=0$ for every other $[\xi]\in \widehat{G}$, but on the other hand
      \begin{align*}
        \|\widehat{{Du}}(\xi_\ell)\|_{\HS}^2&=\sum_{j=1}^{m}\left\|\sum_{i=1}^{n}\sigma_{{D}}(i,j,\xi_\ell)\widehat{u_i}(\xi_\ell)\right\|_{\HS}^2 \\
        &=\sum_{j=1}^{m}\left\|\sum_{i=1}^{n}\sigma_{{D}}(i,j,\xi_\ell)v(i,\xi_\ell)\right\|_{2}^2\\
        &<\langle \xi_\ell\rangle^{-\ell},
    \end{align*}
    for all $\ell\in\N$, while $\widehat{{Du}}(\xi)=0$ for all other $[\xi]\in\widehat{G}$,
    therefore $Du\in C^\infty(G,\C^m)$ and so $D$ is not globally hypoelliptic.
\end{proof}

\end{document}
